% 图 2-2：流水寄存器与反馈通路各占独立走线通道。
\begin{tikzpicture}[font=\footnotesize,
 b/.style={draw=softline,rounded corners=2pt,fill=macblue!55,align=center,text width=2.15cm,inner sep=4pt,minimum height=1.8cm},
 reg/.style={draw=macline,fill=macgray,minimum width=.28cm,minimum height=2.25cm},
 ar/.style={-{Stealth[length=1.4mm]},draw=softline,thick},
 ct/.style={ar,dashed},lab/.style={fill=white,font=\scriptsize,inner sep=2pt}]
\node[b,text width=14cm,minimum height=.8cm,fill=maclav] (ctrl) at (0,0) {\textbf{ctrl.v：集中冒险控制}\\stall / flush[4:0] = \{IF, IF/ID, ID/EX, EX/MEM, MEM/WB\}};
\foreach \x/\s in {-6/IF,-3/ID,0/EX,3/MEM,6/WB}{\node[font=\small\bfseries] at (\x,-1.65) {\s};}
\node[b] (if) at (-6,-3.1) {PC 选择\\I-cache 2 KiB\\两路组相联};
\node[b] (id) at (-3,-3.1) {译码 / 立即数\\寄存器堆\\旁路控制标注};
\node[b] (ex) at (0,-3.1) {旁路 MUX\\ALU / 分支\\乘除 / CSR};
\node[b] (mem) at (3,-3.1) {load / store\\D-cache 2 KiB\\MMIO 绕行};
\node[b] (wb) at (6,-3.1) {结果写回\\GPR 写端口};
\foreach \x/\key/\name/\left/\right in {-4.5/a/{IF/ID}/if/id,-1.5/b/{ID/EX}/id/ex,1.5/c/{EX/MEM}/ex/mem,4.5/d/{MEM/WB}/mem/wb}{
 \node[reg] (r\key) at (\x,-3.1) {};
 \node[rotate=90,font=\scriptsize] at (\x,-3.1) {\name};
 \draw[ar] (\left.east)--(r\key.west);
 \draw[ar] (r\key.east)--(\right.west);
 \draw[ct] (\x,-1.05)--(r\key.north);}
\draw[ct] (ctrl.south)--(0,-1.05)--(-6,-1.05)--(if.north);
\draw[ct] (-4.5,-1.05)--(4.5,-1.05);
\node[b,fill=macblue!55,minimum height=1.3cm] (bp) at (-6,-5.9) {分支预测\\64 项 BTB / BHT};
\draw[ar] (bp.north)--(if.south);
\draw[ct] (ex.south)--(0,-4.65)--(-7.55,-4.65)|-(bp.west);
\draw[ct] (-7.55,-4.65)|-(if.west);
\node[lab,text width=3cm,align=center] at (-4.15,-4.65) {EX 分支重定向\\BTB / BHT 训练};
\node[b,text width=2.6cm,fill=macamber,minimum height=1.3cm] (ise) at (-.35,-5.9) {\textbf{新增}\\\textbf{\mbox{Ascon ISE}}\\320-bit；P8 / P12};
\draw[ar,<->] (ex.south)--(ise.north -| ex.south);
\node[lab] at (0,-4.35) {指令 / 结果};
\draw[ct] (ise.east)--(1.3,-5.9)--(1.3,-.7)--(ctrl.south -| 1.3,0);
\node[lab,anchor=west] at (1.45,-4.65) {ex\_stall\_req};
\draw[ar] (wb.south)--(6,-7.2)--(-2.3,-7.2)--(id.south -| -2.3,0);
\node[lab,anchor=south] at (5.3,-7.16) {GPR 写回};
\draw[ar] (rd.south)--(4.5,-7.85)--(-2.05,-7.85)--(-2.05,-4.25)--(ex.south west);
\draw[ar] (rc.south)--(1.5,-7.85);
\node[lab] at (.4,-7.85) {操作数旁路：EX/MEM $>$ MEM/WB};
\node[b,text width=14cm,minimum height=1.1cm,fill=macgreen] (bus) at (0,-9.2) {I-cache 填充 / D-cache 总线出口 $\rightarrow$ bus\_interconnect\\内部外设；BRAM 经 AXI bridge / master};
\draw[ar] (if.west)--(-7.9,-3.1)--(-7.9,-9.2)--(bus.west);
\draw[ar] (mem.east)--(4.15,-3.1)--(4.15,-8.4)--(3,-8.4)--(bus.north -| 3,0);
\end{tikzpicture}
